% 【本文件写什么】impl 实现层：qemu-loongarch64-virt
% - 定位：QEMU LoongArch64 virt 板级组合。
% - crate 路径：os/components/wateros-driver/driver-impl/impl-qemu-loongarch64-virt/src/
% - 如何实现对应 api-v0 trait：关键类型、状态机、数据结构
% - 算法或硬件交互流程，可用伪代码/流程图
% - 本 impl 启用的 Cargo [features] 与 arch feature，impl-riscv64 等
% - 与其它 impl 变体的差异、切换 feature 时的影响
% - 性能/内存假设与已知限制
% 【事实来源】
% - 上述 crate 源码与 Cargo.toml
% - os/feature-tree.txt 中指向本 impl 的 feature 链

\subsubsection{impl-qemu-loongarch64-virt}

\paragraph{与RISC-V路径差异。}块/网卡不走DTB virtio-mmio扫描，而在\code{init\_after\_boot\_inner}中经PCIe ECAM枚举\code{virtio-blk}/\code{virtio-net}。\code{pci::find\_config\_base}从DTB PCI节点或回退常量解析配置空间基址；\code{probe\_virtio\_blk\_pci}/\code{probe\_virtio\_net\_pci}分别调用块/网\code{impl-virtio-pci}的\code{probe\_first\_from\_ecam}。

\paragraph{引导与RAM。}\code{init\_when\_boot}保存DTB并\code{uart::init\_early\_default\_uart}初始化\code{0x1FE0\_01E0} NS16550。\code{physical\_ram\_end\_exclusive}优先DTB\code{memory@*}，失败回退\code{0xb000\_0000}以匹配QEMU\code{virt -m 1G}高段布局；内核自 \code{0x9000\_0000} 启动，避免把MMIO空洞当作RAM。

\paragraph{注册与devfs。}块设备成功时可选\code{BlockCacheManager::wrap}并\code{register\_block\_device}，记录\code{VirtioPciProbeInfo}到\code{VIRTIO\_BLK\_PCI}。网卡\code{register\_network\_device}并记录MAC与BDF。字符路径调用\code{uart::register\_default\_uart}，全局 \code{QemuLoongArch64Uart16550}，与 \code{register\_builtin\_character\_devices}。末尾\code{devfs\_impl::refresh}暴露\code{/dev/vblk*}等节点。

\begin{rustcode}
pub fn init_when_boot(dtb_pa: usize) {
    DTB_BASE_ADDR.store(dtb_pa, Ordering::Release);
    uart::init_early_default_uart();
}

pub fn physical_ram_end_exclusive() -> usize {
    // DTB memory@* 或回退 0xb000_0000
    0xb000_0000
}
\end{rustcode}

\paragraph{导出接口。}除与RISC-V相同的\code{init\_after\_boot}/\code{physical\_ram\_end\_exclusive}/\code{test}外，\code{uart::with\_default\_uart}供\code{klog}与早期打印在已初始化UART上执行闭包。重复\code{init\_after\_boot}同样幂等忽略。

\paragraph{限制。}仅扫描PCI bus0首个块/网设备；DTB扫描不用于virtio绑定，\code{scan\_device\_info} 非本路径主流程。中断未启用，IO与收包依赖轮询。多实例devfs与syscall策略仍以首设备为主。
